\subsection{3 марта. Сопряжённый оператор}

\begin{definition}
Пусть $E_1, E_2$~--- линейные нормированные пространства,
$A \in \mathcal{L}(E_1, E_2)$. \textit{Сопряжённым} к оператору $A$ называется
$A^*: E_2^* \to E_1^*$, определённый по правилу
$(A^* g)(x) = g(A x) \ \forall g \in E_2^*$.  
% TODO дописать и вставить картинки
\end{definition}



\begin{theorem}[9.1]
    Пусть $E_1, E_2$~--- линейные нормированные пространства. Тогда $\forall A \in \mathcal{L}(E_1, E_2) \; \exists ! A^*$, причем $\| A^* \| =\| A\|$
\end{theorem}
\begin{proof}
    Существование и единственность очевидны из определения
    Ограниченность оператора $A^*$:

    \[
    |(A^*g)(x)| = |g(Ax)| \leqslant \|g\| \cdot \|Ax\| \leqslant
    \|g\| \cdot \|A\| \cdot \|x\|
    \]
    
    \[
    \|A^* g\| \leqslant \|g\| \cdot \|A\| \Rightarrow
    \|A^*\| \leqslant \|A\|
    \]
    % кухня Хана-Банаха
    С другой стороны, $\|A\| \leqslant \|A^*\|$, так как по следствию 4 теоремы Хана-Банаха
    \[ \| Ax\| = \sup_{\|g\|_{E_2^*} = 1} |g(Ax)| = 
    \sup_{\|g\|=1} |(A^*g)(x)| \leqslant \|A^*\| \|x\|\]
\end{proof}

\begin{theorem}[9.3 (б/д)]
Пусть $E_1, E_2$~--- линейные нормированные пространства. Тогда
$\exists A^{-1} \in \mathcal{L}(E_2, E_1)$ тогда и только тогда, когда
$\exists (A^*)^{-1} \in \mathcal{L}(E_1^*, E_2^*)$.
При этом $(A^*)^{-1} = (A^{-1})^*$.
\end{theorem}

Ситуация: $E_1 = H_1, E_2 = H_2$~--- гильбертовы пространства.
В философии Рисса-Фреше:

\[ (x, A^*y)_{H_1}=(A^* g)(x)=g(Ax)=(Ax,y)_{H_2}\]

Получается, что эрмитово сопряжение является частным случаем сопряжения при $H_1=H_2$.

$H(\C), A \in \mathcal{L}(H)$

\begin{exercise}
Доказать, что сопряжённый оператор (по Эрмиту) всегда существует и единственен. Доказать свойства:

\begin{enumerate}
    \item $I^* = I$
    \item $(\alpha A + \beta B)^* = \overline{\alpha} A^* + \overline{\beta} B^*$
    \item $(AB)^* = B^* A^*$
\end{enumerate}

\end{exercise}

\begin{lemmanote}
    Как искать $A^*$? Общего решения нет, но берем и мучаем $(Ax, y) = \ldots$ до тех пор, пока оно не предствится в виде $(x, \zeta)$, 
    где $\zeta$~--- неприятная буква греческого алфавита.
\end{lemmanote}

\begin{example}
\begin{enumerate}
    \item $\R^n_2, \; A \sim (a_{ij}), \; A^* \sim (a_{ji})$
    \item $L_2[0,1]$, оператор Вольтерра $(V f)(x) = \int_0^x f(t) dt$.
    
    Как найти $V^*$? (Можно сразу понять, что $\|V^*\| = 1$).
\end{enumerate}
\end{example}

\begin{theorem}[9.2]\label{9.2}
    $H(\K)$ - гильбертово пространство, $A \in \mathcal{L}(H)$.. 
    Тогда $H = \overline{\im A} \oplus_{\perp} \ker A^*$
\end{theorem}

\begin{note}
$H = \overline{\im A^*} \oplus \ker A$
\end{note}

\begin{proof}
    \begin{enumerate}
        \item Простая часть: $(\im A)^\perp = \ker A^*$. Почему?

        \[ y \in (\im A)^\perp \sim \forall x \in \im A \; (x,y)=0 \sim \forall z \in H \; (Az, y)=0\]

        Тогда $\forall z \in H \; (z, A^*y)=0 \sim A^*y=0 \sim y \in \ker A^*$
        % здесь история про греков и гору (и рыбку!!!)
        \item <<Сложная>> часть:
        Знаем, что для любого множества $S$ верно, что $S^{\perp} = \overline{S}^{\perp}$.

        $(Im A)^{\perp} = \ker A^*$, где $\im A$~--- линейное многообразие,
        $\left( \overline{\im A} \right)^{\perp}$~--- подпространство.

        По теореме 4.3 (о проекции) получаем 
        \[ \overline{(\im A)} \oplus \left( \overline{\im A} \right)^{\perp}=\overline{(\im A)} \oplus \ker A^* = H \]
    \end{enumerate}
\end{proof}

\begin{theorem}[Банаха об обратном операторе (8.4) в случае гильбертовых пространств]\label{8.4}
     Пусть $H(\K)$~--- гильбртово простраство, $A \in \mathcal{L}(H)$,
     $A$~--- биекция $H$ на $H$. Тогда $A^{-1} \in \mathcal{L}(H)$.
\end{theorem}

\begin{proof}  % доказательство, которое невозможно забыть
    \begin{itemize}
        \item[0.] Докажем, что из условия теоремы следует существование $(A^*)^{-1} \in \L(H)$  % шило на мыло
        \item[1.] Рассмотрим формулы в теореме \nameref{9.2}
        Так как $\ker A = \{0\}$, $\im A = H$, то
        $\ker A^* = \{0\}$ (это означает, что
        $A^*$~--- взаимно-однозначный) и $\overline{\im A^*} = H$.
        Есть теорема \nameref{8.1}, по которой достаточно показать, что
        $A^*$~--- ограниченный снизу, то есть 
        $\exists m > 0:\ \forall x \ \|A^* x\| \geqslant m \|x\|$.
        \item[2.] \textbf{Лемма:} $E$~--- банахово пространство,
        $A \in \L(E)$, $A$~--- ограниченный снизу. 
        Тогда $\overline{\im A} = \im A$.
        
        \textbf{Доказательство леммы:} Пусть $y_n \in \im A, \forall n$, $y_n \to y$. Следует ли из этого, что $y \in \im A$?
        $y_n = A z_n, z_n \in E$. $\{y_n\}$ сходится, а значит она фундаментальна.
    
        \[ 0 \leftarrow \|y_{n+p}-y_n\|=\|Az_{n+p}-Az_n\| \geq m \cdot \|z_{n+p} - z_n \| \to 0 \] 

        Итак, $\{z_n\}$ - фундаментальна в $E$ - банахово, откуда $\exists \lim z_n = z$. Следовательно, $y_n = A z_n \to A z = y$. То есть $y \in \im A \; \blacksquare$
        \item[3.] Осталось показать, что $A^*$~--- ограниченный снизу, то есть $\exists m > 0 : \|A^* x\| \leqslant m \|x\| \ \forall x \in H$.

        Вспомним, что $\ker A^* = \{0\}$. Если неравенство верно для $x$, то оно также верно для $\alpha x$ ($\alpha \neq 0$).

        \textbf{Трюк:} рассмотрим множество 
        $S = \{x \in H \: | \: \|A^*x\| = 1\}$. Следовательно, достаточно показать, что $\exists m > 0 : \forall x \in S \; 1 \geqslant m \|x\| \sim \|x\| \leqslant \frac{1}{m}$. То есть надо доказать, что $S$~--- ограниченное множество.
        % Греция -- штука странная. Куда плавал Одиссей?
        Вспомним теорему Хана из лекции 2.
        % Будете себя плохо вести -- докажем её ещё раз
        Покажем, что $S$~--- слабо ограниченное, то есть
        \[
        \forall y \in H \exists K_y : |(x, y)| \leqslant K_y \; \forall x \in S
        \]

        % главное место в доказательстве
        Используем, что $A$~--- биекция $H$ на $H$: $\forall y \in H \; \exists z \in H \; Az = y$

        \[ |(x, y)| = |(x, Az)| = |A^*x, z)| \leq \underbrace{\|A^* x\|}_{=1} \cdot \underbrace{\|z\|}_{=K_y}\]
    \end{itemize}
    
    Из условия теоремы следует, что $\exists (A^*)^{-1} \in \L(H)$.
    Тогда $A^*$ удовлетворяет условиям нашей теоремы.
    Следовательно, $(\underbrace{A^{**}}_{=A})^{-1} \in \L(H)$.
\end{proof}